\section*{अध्याय \ref{ch:b1:quantum-introduction} --- क्वांटम भौतिकी का परिचय}

\begin{solution}{exo:b1:quantum-introduction:1}
$h\nu$: \qty{100}{MHz}: $6.6 \times 10^{-26}$~जूल $= \qty{4.1e-7}{eV}$; \qty{500}{nm}:
$1240/500 = \qty{2.5}{eV}$; \qty{0.10}{nm}: \qty{12.4}{keV}; \qty{1}{MeV}: $\lambda =
1240/10^6 = \qty{1.2e-3}{nm} = \qty{1.2}{pm}$। लेज़र: $1240/633 = \qty{1.96}{eV} =
\qty{3.1e-19}{J}$, अतः प्रति सेकंड $10^{-3}/3.1 \times 10^{-19} = 3.2 \times 10^{15}$
फ़ोटॉन।
\end{solution}

\begin{solution}{exo:b1:quantum-introduction:2}
$\lambda_0 = 1240/2.3 = \qty{540}{nm}$। \qty{400}{nm} पर, $h\nu = \qty{3.1}{eV}$:
$E_{k,\max} = \qty{0.8}{eV}$, $V_s = \qty{0.8}{V}$। तीव्रता दुगुनी करने पर
धारा दुगुनी हो जाती है और $V_s$ वही रहता है। और \qty{600}{nm} पर
(\qty{2.07}{eV} $< W$): कुछ भी नहीं, किसी भी तीव्रता पर।
\end{solution}

\begin{solution}{exo:b1:quantum-introduction:3}
इलेक्ट्रॉन: $1.23/\sqrt{100} = \qty{0.12}{nm}$; और \qty{1}{keV} पर: \qty{0.039}{nm}।
न्यूट्रॉन: $E = \qty{4.0e-21}{J}$, $p = \sqrt{2m_nE} = \qty{3.7e-24}{kg.m/s}$,
$\lambda = h/p = \qty{0.18}{nm}$। गेंद: $6.6 \times 10^{-34}/10^{-3} = \qty{7e-31}{m}$।
पहले तीनों की $\lambda$ परमाणुओं की दूरी की कोटि की है और वे विवर्तित होते
हैं (इलेक्ट्रॉन, न्यूट्रॉन और X-किरण रवाविज्ञान); गेंद कभी नहीं।
\end{solution}

\begin{solution}{exo:b1:quantum-introduction:4}
$E_1 = h^2/8m_eL^2 = 4.39 \times 10^{-67}/(8 \times 9.11 \times 10^{-31} \times 10^{-18}) =
\qty{6.0e-20}{J} = \qty{0.38}{eV}$, $E_2 = 4E_1 = \qty{1.5}{eV}$; फ़ोटॉन
\qty{1.13}{eV}, $\lambda = 1240/1.13 = \qty{1.1}{\micro m}$ (निकट अवरक्त)।
\qty{5}{fm} में प्रोटॉन: $4.39 \times 10^{-67}/(8 \times 1.67 \times 10^{-27} \times 25 \times 10^{-30}) =
\qty{1.3e-12}{J} = \qty{8}{MeV}$ --- यानी नाभिकीय ऊर्जाओं का पैमाना, केवल
बंदी से।
\end{solution}

\begin{solution}{exo:b1:quantum-introduction:5}
ढाल $(2.05 - 0.40)/(4.0 \times 10^{14}) = \qty{4.1e-15}{V.s}$, अतः $h = e \times 4.1 \times
10^{-15} = \qty{6.6e-34}{J.s}$। $W = h\nu - eV_s = 2.48 - 0.40 = \qty{2.1}{eV}$,
और देहली $\lambda_0 = 1240/2.1 = \qty{600}{nm}$: यानी सीज़ियम।
\end{solution}

\begin{solution}{exo:b1:quantum-introduction:6}
$\Delta\lambda = \qty{2.43}{pm}$: अतः $\lambda' = \qty{12.4}{pm}$। पहले: $1240/0.010 =
\qty{124}{keV}$; बाद में: $1240/0.0124 = \qty{100}{keV}$; और इलेक्ट्रॉन
\qty{24}{keV} ले जाता है। प्रकाश के लिए, $\Delta\lambda/\lambda = 2.4 \times 10^{-12}/5 \times 10^{-7} = 5 \times
10^{-6}$: यानी 1923 में अमाप्य --- X-किरणों ने उस विस्थापन को $25\%$ का प्रभाव बना दिया।
\end{solution}

\begin{solution}{exo:b1:quantum-introduction:7}
$p = \sqrt{2 \times 9.11 \times 10^{-31} \times 5 \times 10^4 \times 1.6 \times 10^{-19}} =
\qty{1.2e-22}{kg.m/s}$, $\lambda = \qty{5.5}{pm}$ (आपेक्षिकता उसे
\qty{5.4}{pm} तक सुधार देती है); और फ्रिंजों की दूरी $\lambda D/d = 5.5 \times 10^{-12} \times 0.35/2 \times 10^{-6}
= \qty{1}{\micro m}$: यानी आँख के लिए अदृश्य, इसीलिए \omterm{def:b1:optical-instruments:G}{आवर्धक} \omterm{def:b1:mirrors-thin-lenses:lens}{लेंस}। दस
इलेक्ट्रॉन: दस बिंदु, देखने में यादृच्छिक; और $10^5$: फ्रिंजें। हर
इलेक्ट्रॉन किसी एक बिंदु पर आ गिरता है; पर गिरने की \emph{प्रायिकता}
दो-झिर्री तरंग का पालन करती है --- यानी अकेले इलेक्ट्रॉन की तरंग दोनों
झिर्रियों से गुज़रती है।
\end{solution}

\begin{solution}{exo:b1:quantum-introduction:8}
परमाणु: $\Delta p \geq \hbar/2\Delta x = \qty{5.3e-25}{kg.m/s}$, $E \sim \Delta p^2/2m_e =
\qty{1.5e-19}{J} \approx \qty{1}{eV}$: यानी सही पैमाना। नाभिक: $\Delta p \geq
\qty{1.1e-20}{kg.m/s}$; और $\Delta p^2/2m_e$ \qty{400}{MeV} निकलता, जो
$m_ec^2$ से कहीं ऊपर है, अतः $E \approx \Delta p\,c = \qty{3e-12}{J} = \qty{20}{MeV}$
लीजिए: नाभिक में कोई भी बल इतने ऊर्जावान इलेक्ट्रॉन को थाम नहीं सकता ---
यानी नाभिकों में इलेक्ट्रॉन होते ही नहीं (बीटा इलेक्ट्रॉन उत्सर्जन के समय ही
बनते हैं)। कण: $\Delta v \geq
\hbar/2m\Delta x = 1.05 \times 10^{-34}/(2 \times 10^{-9} \times 10^{-6}) = \qty{5e-20}{m/s}$:
यानी अप्रासंगिक।
\end{solution}

\begin{solution}{exo:b1:quantum-introduction:9}
(क) $m_ev^2/r = e^2/4\pi\varepsilon_0r^2$: $v = \sqrt{e^2/4\pi\varepsilon_0m_er}$, $E =
\tfrac12m_ev^2 - e^2/4\pi\varepsilon_0r = -e^2/8\pi\varepsilon_0r$। (ख) $m_evr = n\hbar$ और
$m_ev^2r = e^2/4\pi\varepsilon_0$: पहले के वर्ग को दूसरे से भाग दीजिए:
$m_er = n^2\hbar^2 \cdot 4\pi\varepsilon_0/e^2$, अतः $r_n = n^2a_0$, और $E_n = -e^2/8\pi\varepsilon_0n^2a_0
= E_1/n^2$, जहाँ $a_0 = \qty{52.9}{pm}$, $E_1 = \qty{-13.6}{eV}$। (ग) $v_1 =
\hbar/m_ea_0 = \qty{2.19e6}{m/s}$, $v_1/c = 1/137$। (घ) $2 \to 1$: $\tfrac34 \times
13.6 = \qty{10.2}{eV}$, \qty{122}{nm}; और $3 \to 2$: $(\tfrac14 - \tfrac19) \times 13.6 =
\qty{1.89}{eV}$, \qty{656}{nm}। (ङ) $\lambda_n = h/m_ev_n = nh/m_ev_1 = 2\pi na_0$
(क्योंकि $m_ev_1a_0 = \hbar$), और $2\pi r_n = 2\pi n^2a_0 = n\lambda_n$।
\end{solution}

\begin{solution}{exo:b1:quantum-introduction:10}
(क) $e^2$ की जगह $Ze^2$ रखिए: $a_0 \to a_0/Z$, $E_1 \to Z^2E_1$। He$^+$:
\qty{54.4}{eV}। U$^{91+}$: $92^2 \times 13.6 = \qty{115}{keV}$ --- पर $v_1 = Z\alpha c
= 0.67c$: अतः वहाँ आपेक्षिकता-रहित प्रतिरूप केवल संकेत भर देता है। (ख) $a_0/207
= \qty{256}{fm}$, $E_1 = 207 \times 13.6 = \qty{2.8}{keV}$; और लाइमन $\alpha$:
\qty{2.1}{keV}, $\lambda = \qty{0.59}{nm}$ (X-किरणें)। म्यूऑन की कक्षा केवल
$300$ प्रोटॉन-त्रिज्याओं की है: अतः उसके स्तर प्रोटॉन के आकार के साथ मापने
लायक खिसक जाते हैं, और इसी से प्रोटॉन की त्रिज्या फिर से नापी गई (2010)।
(ग) दोनों कण अपने उभयनिष्ठ केंद्र का चक्कर लगाते हैं: \omterm{prop:b1:systems-of-points:twobody}{समानीत द्रव्यमान}
$m_e/2$ हर ऊर्जा को आधा कर देता है: $E_1 = \qty{-6.8}{eV}$, और लाइमन $\alpha$ \qty{5.1}{eV} पर, \qty{243}{nm}।
\end{solution}

\begin{solution}{exo:b1:quantum-introduction:11}
(क) $h^2/8m_e = \qty{6.02e-38}{J.m^2}$। $R = \qty{2}{nm}$: $E_e = 6.02 \times 10^{-38}/
(0.13 \times 4 \times 10^{-18}) = \qty{0.72}{eV}$, $E_h = \qty{0.21}{eV}$, कुल $1.74 +
0.93 = \qty{2.67}{eV}$, \qty{464}{nm}: यानी नीला। \qty{3}{nm}: बंदी-ऊर्जा $\times4/9$,
\qty{2.15}{eV}, \qty{577}{nm}: पीला। \qty{4}{nm}: $\times1/4$, \qty{1.97}{eV},
\qty{629}{nm}: लाल। (ख) \qty{520}{nm} यानी \qty{2.38}{eV}: बंदी-ऊर्जा \qty{0.65}{eV}
$= 0.93 \times (2/R)^2$, अतः $R = \qty{2.4}{nm}$। (ग) कोई बंदी नहीं: अकेला $E_g$,
\qty{713}{nm}। (घ) $\delta E = -2 \times 0.1 \times 0.41 = \qty{-0.08}{eV}$, यानी
\qty{2.15}{eV} का $4\%$: लगभग \qty{20}{nm} का फैलाव --- शुद्ध रंगों के लिए
एक-से आकार के बैच चाहिए।
\end{solution}

\begin{solution}{exo:b1:quantum-introduction:12}
(क) $E(\Delta x) = \hbar^2/8m\Delta x^2 + \tfrac12m\omega^2\Delta x^2$; और अवकलज शून्य
$\Delta x^2 = \hbar/2m\omega$ पर, जहाँ दोनों पद $\hbar\omega/4$ के बराबर हैं: अतः $E_{\min} =
\hbar\omega/2$। (ख) $\tfrac12 \times 1.055 \times 10^{-34} \times 8.3 \times 10^{14} = \qty{4.4e-20}{J}
= \qty{0.27}{eV}$; और $T \sim \hbar\omega/k_B = \qty{6300}{K}$ --- यानी कमरे के ताप
पर कंपन अपनी \omterm{prop:b1:quantum-introduction:box}{मूल अवस्था} में जमा रहता है। (ग) $\Delta p = \hbar/2\Delta x =
\qty{1.05e-24}{kg.m/s}$, $E = \Delta p^2/2m_{\mathrm{He}} = 1.1 \times 10^{-48}/1.33 \times
10^{-26} = \qty{8e-23}{J} = \qty{0.5}{meV}$: यानी बंधन का आधा --- परमाणु इतने
स्थिर बैठ ही नहीं सकते कि रवा बना सकें; अतः हीलियम परम शून्य तक द्रव बना
रहता है, जब तक उसे दबाया न जाए (\qty{25}{bar})। (घ) $E_1 = \pi^2\hbar^2/2mL^2$, जबकि
$\hbar^2/2mL^2$: यानी दस गुना बड़ा, और होना भी चाहिए --- सीमा तो कोई निचली
हद है, और डिब्बे वाली अवस्था का $\Delta x \approx 0.18L$ है, $L/2$ नहीं।
\end{solution}

\begin{solution}{pb:b1:quantum-introduction:1}
\textbf{1.} प्रकाश कैथोड से इलेक्ट्रॉनों को भाँति-भाँति की ऊर्जाओं के साथ
मुक्त करता है; और $-V_s$ पर रखा ऐनोड उन्हें प्रतिकर्षित करता है, अतः केवल
वही पहुँचते हैं जिनकी $E_k > eV_s$ हो। धारा तब लुप्त होती है जब $eV_s = E_{k,\max}$।

\textbf{2.} $eV_s = h\nu - W$: ढाल $h/e$, अंतःखंड $-W/e$;
और देहली $\nu_0 = W/h$ पर $V_s = 0$।

\textbf{3.} चरम बिंदु: $(2.36 - 0.30)/(5.0 \times 10^{14}) = \qty{4.12e-15}{V.s}$,
अतः $h = 1.602 \times 10^{-19} \times 4.12 \times 10^{-15} = \qty{6.60e-34}{J.s}$ (पाँचों
पर फ़िट करने से $6.62$ मिलता है): यानी स्वीकृत $6.626 \times 10^{-34}$ से
$0.5\%$ के भीतर।

\textbf{4.} $W = h\nu - eV_s = 2.47 - 0.30 = \qty{2.2}{eV}$; $\nu_0 = W/h =
\qty{5.3e14}{Hz}$, $\lambda_0 = 1240/2.2 = \qty{560}{nm}$। कोई लाल सूचक
(\qty{650}{nm}, \qty{1.9}{eV}) कुछ भी नहीं निकालता।

\textbf{5.} $V_s$ वही, और धारा दुगुनी: यानी दुगुने फ़ोटॉन, हर एक उसी ऊर्जा
का। कोई तरंग तो अधिक तीव्रता पर हर इलेक्ट्रॉन को अधिक ऊर्जा देती, और देहली
आवृत्ति में नहीं बल्कि तीव्रता में होती।

\textbf{6.} परमाणु पर $10^{-6} \times 10^{-20} = \qty{1e-26}{W}$; और \qty{2.2}{eV} $=
\qty{3.5e-19}{J}$ जमा होने में $3.5 \times 10^7$~सेकंड --- यानी एक बरस। जबकि
उत्सर्जन नैनोसेकंडों में देखा जाता है: ऊर्जा पूरे-पूरे क्वांटमों में आती है।

\textbf{7.} $m_ev^2/r = e^2/4\pi\varepsilon_0r^2$: $v = \sqrt{e^2/4\pi\varepsilon_0m_er}$; $E =
\tfrac12m_ev^2 - e^2/4\pi\varepsilon_0r = -e^2/8\pi\varepsilon_0r$।

\textbf{8.} $(m_evr)^2 = n^2\hbar^2$ और $m_ev^2r = e^2/4\pi\varepsilon_0$: अतः $r_n = n^2 \cdot
4\pi\varepsilon_0\hbar^2/m_ee^2 = n^2a_0$, और $a_0 = (1.055 \times 10^{-34})^2/(8.99 \times 10^9
\times 9.11 \times 10^{-31} \times 2.57 \times 10^{-38}) = \qty{52.9}{pm}$।

\textbf{9.} $E_n = -e^2/8\pi\varepsilon_0r_n = -m_ee^4/2(4\pi\varepsilon_0)^2\hbar^2n^2$; $E_1 =
-2.31 \times 10^{-28}/(2 \times 5.29 \times 10^{-11}) = \qty{-2.18e-18}{J} = \qty{-13.6}{eV}$।

\textbf{10.} $v_1 = \hbar/m_ea_0 = \qty{2.19e6}{m/s}$; $v_1/c = 7.3 \times 10^{-3} =
1/137$।

\textbf{11.} $2 \to 1$: \qty{10.2}{eV}, \qty{122}{nm} (पराबैंगनी); $3 \to 2$:
\qty{1.89}{eV}, \qty{656}{nm} (लाल); $4 \to 2$: \qty{2.55}{eV}, \qty{486}{nm}
(नीला-हरा): और दृश्य वही बामर रेखाएँ हैं।

\textbf{12.} $hc/\lambda = |E_1|(1/n^2 - 1/p^2)$: अतः $R_H = |E_1|/hc = 13.6/1240 =
\qty{1.097e-2}{nm^{-1}} = \qty{1.097e7}{m^{-1}}$। बामर सीमा ($p \to \infty$):
$\lambda = 4/R_H = \qty{365}{nm}$।

\textbf{13.} \qty{13.6}{eV}; $\lambda = 1240/13.6 = \qty{91}{nm}$।

\textbf{14.} $v_n = v_1/n$, $\lambda_n = h/m_ev_n = nh/m_ev_1 = 2\pi na_0$ (क्योंकि $m_ev_1a_0
= \hbar$); और $2\pi r_n = 2\pi n^2a_0 = n\lambda_n$: यानी कक्षा $n$ पूरी-पूरी
तरंगें ढोती है --- बोर का नियम कोई अप्रगामी-तरंग नियम ही है।

\textbf{15.} ग़लत: इलेक्ट्रॉन की कोई कक्षा होती ही नहीं (वह विकिरण करता; और
\omterm{prop:b1:quantum-introduction:box}{मूल अवस्था} का \omterm{def:b1:angular-momentum:L}{कोणीय संवेग} शून्य है), और यह प्रतिरूप दो इलेक्ट्रॉनों वाले
किसी भी परमाणु के लिए विफल हो जाता है। ठीक: हाइड्रोजन के स्तर, और इसलिए
उसके वर्णक्रम की हर रेखा। कक्षा की जगह कोई \omterm{def:b1:quantum-introduction:wavefunction}{तरंग-फलन} ले लेता है, यानी
$\sim n^2a_0$ आकार का कोई प्रायिकता-मेघ।

\textbf{16.} अप्रगामी तरंगें: $L = n\lambda/2$, $p_n = h/\lambda_n = nh/2L$, $E_n =
p_n^2/2m = n^2h^2/8mL^2$। और गोले वाला सूत्र वही है, जहाँ $L \to R$:
\omterm{prop:b1:quantum-introduction:box}{मूल अवस्था} त्रिज्या में आधी तरंगदैर्घ्य बिठाती है।

\textbf{17.} $E_e = 6.02 \times 10^{-38}/(0.13 \times 4 \times 10^{-18}) = \qty{0.72}{eV}$;
$E_h = 0.13/0.45 \times 0.72 = \qty{0.21}{eV}$।

\textbf{18.} $R = \qty{2}{nm}$: $1.74 + 0.93 = \qty{2.67}{eV}$, \qty{464}{nm}, नीला;
\qty{3}{nm}: \qty{2.15}{eV}, \qty{577}{nm}, पीला; \qty{4}{nm}: \qty{1.97}{eV},
\qty{629}{nm}, लाल।

\textbf{19.} बंदी-ऊर्जा $1/R^2$ की तरह बढ़ती है: छोटा डिब्बा स्तरों को ऊपर
उठा देता है, फ़ोटॉन अधिक ऊर्जावान होता है, और प्रकाश अधिक नीला।

\textbf{20.} $\Delta p \sim \hbar/R = \qty{5.3e-26}{kg.m/s}$, $E \sim \Delta p^2/2m_e^\ast =
\qty{1.2e-20}{J} = \qty{0.07}{eV}$: यानी सही कोटि (यथार्थ गुणक
$\pi^2/2$ और मोटा-मोटी लिया गया $\Delta p$ मिलकर दस के गुणक का हिसाब कर देते हैं)।

\textbf{21.} $1240/577 = \qty{2.15}{eV} = \qty{3.4e-19}{J}$: अतः प्रति सेकंड $10^{-9}/3.4 \times 10^{-19}
= 2.9 \times 10^9$ फ़ोटॉन।

\textbf{22.} प्रति सेकंड $10^{-3}/(1.96 \times 1.6 \times 10^{-19}) = 3.2 \times 10^{15}$।

\textbf{23.} \qty{0.1}{s} में $100 \times 2.48 \times 1.6 \times 10^{-19} = \qty{4e-17}{J}$:
यानी \qty{4e-16}{W}।

\textbf{24.} $h\nu = \qty{6.6e-26}{J}$: प्रति सेकंड $1.5 \times 10^{13}$ फ़ोटॉन ---
और हर एक एंटीना की तापीय ऊर्जा $k_BT$ से $10^5$ गुना छोटा: अतः कोई आशा
नहीं, और कोई ज़रूरत भी नहीं; वहाँ तरंग-चित्र ही यथार्थ है।

\textbf{25.} परमाणुओं का आकार, $a_0 = \qty{53}{pm}$; उनकी ऊर्जा का पैमाना,
\qty{13.6}{eV} (और इसीलिए \omterm{def:b1:charged-particles:eV}{इलेक्ट्रॉन-वोल्ट} वाले फ़ोटॉन, रसायन, दृश्य
वर्णक्रम); और उनके इलेक्ट्रॉनों की चाल, $\alpha c \approx \qty{2e6}{m/s}$ ---
ये तीनों $h$, $e$ और $m_e$ से (तथा $\varepsilon_0$ से), बिना कुछ भी समायोजित किए।
\end{solution}
